\documentclass[10pt,letterpaper]{article}

% ---------- Page and typography ----------
\usepackage[
  top=0.40in,
  bottom=0.40in,
  left=0.60in,
  right=0.60in
]{geometry}
\usepackage[T1]{fontenc}
\usepackage{lmodern}
\usepackage{microtype}
\usepackage{enumitem}
\usepackage{tabularx}
\usepackage{titlesec}
\usepackage{xcolor}
\usepackage[hidelinks]{hyperref}
\usepackage{iftex}

% Make the generated PDF text searchable and ATS-friendly when using pdfLaTeX.
\ifPDFTeX
  \input{glyphtounicode}
  \pdfgentounicode=1
\fi

\definecolor{navy}{HTML}{17324D}
\definecolor{blue}{HTML}{247BA0}
\definecolor{charcoal}{HTML}{25313C}
\definecolor{muted}{HTML}{566573}

\hypersetup{
  colorlinks=true,
  urlcolor=navy,
  linkcolor=navy,
  pdftitle={Maya Chen - Biomedical Engineer Resume},
  pdfauthor={Maya Chen},
  pdfsubject={Biomedical engineering resume},
  pdfkeywords={biomedical engineer, medical devices, verification, validation, ISO 13485, FDA, design controls}
}

\pagestyle{empty}
\setlength{\parindent}{0pt}
\setlength{\parskip}{0pt}
\setlist[itemize]{
  leftmargin=1.15em,
  label=\textcolor{blue}{\small\textbullet},
  itemsep=0.7pt,
  topsep=1.8pt,
  parsep=0pt,
  partopsep=0pt
}
\urlstyle{same}

\titleformat{\section}
  {\fontsize{11.3}{12.2}\selectfont\bfseries\color{navy}}
  {}{0pt}{}
  [\vspace{-0.35em}\textcolor{blue}{\rule{\textwidth}{0.7pt}}]
\titlespacing*{\section}{0pt}{4.5pt}{2.7pt}

% ---------- Reusable resume commands ----------
\newcommand{\contactsep}{\hspace{0.55em}\textcolor{blue}{\textbar}\hspace{0.55em}}

\newcommand{\role}[4]{%
  \begin{tabularx}{\textwidth}{@{}Xr@{}}
    \textbf{\color{charcoal}#1} & \textbf{\color{charcoal}#2} \\
    \textit{\color{muted}#3} & \textit{\color{muted}#4}
  \end{tabularx}\vspace{-0.25em}
}

\newcommand{\project}[2]{%
  \textbf{\color{charcoal}#1} \hfill \textit{\color{muted}#2}\par
}

\newcommand{\skillrow}[2]{%
  \textbf{\color{charcoal}#1:} #2\par
}

\begin{document}
\fontsize{9.35}{10.7}\selectfont

% ---------- Header ----------
\begin{center}
  {\fontsize{23}{25}\selectfont\bfseries\color{navy} MAYA CHEN}\par
  \vspace{1pt}
  {\fontsize{11.5}{12.5}\selectfont\bfseries\color{blue} BIOMEDICAL ENGINEER}\par
  \vspace{4pt}
  \small
  Boston, MA
  \contactsep
  \href{tel:+16175550184}{(617) 555-0184}
  \contactsep
  \href{mailto:maya.chen@protonmail.com}{maya.chen@protonmail.com}
  \contactsep
  \href{https://www.linkedin.com/in/maya-chen-bme}{linkedin.com/in/maya-chen-bme}
\end{center}

\vspace{-2pt}

% ---------- Professional summary ----------
\section*{PROFESSIONAL SUMMARY}
Biomedical engineer with 6+ years of experience developing Class II electromechanical and wearable medical devices from concept through design transfer. Combines hands-on prototyping and test-method development with rigorous design controls, risk management, and cross-functional leadership. Delivered three product launches, reduced verification cycle time by 28\%, and supported successful FDA 510(k) submissions and ISO 13485 audits.

% ---------- Experience ----------
\section*{PROFESSIONAL EXPERIENCE}

\role{Senior Biomedical Engineer}{Jan 2022 -- Present}{Helix Medical Systems}{Cambridge, MA}
\begin{itemize}
  \item Lead systems engineering for a portable infusion platform, translating 120+ clinical and user needs into traceable inputs, subsystem requirements, and verification protocols for a Class II device.
  \item Designed benchtop fixtures and automated Python data pipelines for flow accuracy, occlusion, and battery testing, cutting execution and analysis time by 28\% across 1,500+ test runs.
  \item Drove ISO 14971 risk activities, including hazard analysis, FMEAs, and risk-control verification; closed 34 high-priority risks before design freeze with no unresolved unacceptable residual risk.
  \item Coordinated firmware, mechanical, quality, clinical, and manufacturing teams through design reviews and transfer, enabling pilot production at 96\% first-pass yield and an on-time 510(k) submission.
  \item Mentored four engineers and standardized verification templates, calibration practices, and statistical acceptance criteria across two product programs.
\end{itemize}

\role{Biomedical Engineer II}{Jul 2019 -- Dec 2021}{VitaSense Technologies}{Waltham, MA}
\begin{itemize}
  \item Developed and verified a wearable cardiorespiratory monitor using ECG, PPG, and inertial sensors; improved motion-artifact rejection by 22\% through signal characterization and algorithm tuning.
  \item Created test methods in MATLAB and LabVIEW, performed Gage R\&R and correlation studies, and authored 40+ protocols and reports under ISO 13485 design controls.
  \item Investigated field and manufacturing failures using fault-tree analysis, 5 Whys, and DOE; corrective actions lowered sensor-related nonconformances by 31\% in six months.
  \item Partnered with clinicians on formative usability studies and converted findings into ergonomic and interface changes that increased successful task completion from 82\% to 97\%.
\end{itemize}

\role{Biomedical Engineering Co-op}{Jan 2018 -- Jun 2019}{Northeastern University Center for Translational NeuroImaging}{Boston, MA}
\begin{itemize}
  \item Built a synchronized EMG and force-sensing acquisition system for upper-limb rehabilitation research, reducing experimental setup time by 35\% and improving signal repeatability.
  \item Designed SolidWorks enclosures and 3D-printed fixtures, validated sensor calibration, and maintained IRB-compliant test documentation for studies involving 60+ participants.
\end{itemize}

% ---------- Projects ----------
\section*{SELECTED ENGINEERING PROJECTS}

\project{Low-Cost Prosthetic Hand Controller}{Embedded Systems, Signal Processing}
Designed an eight-channel EMG acquisition board and real-time gesture classifier on an STM32 microcontroller; achieved 93\% classification accuracy across six grasp patterns with sub-100 ms response time.

\vspace{2pt}
\project{MRI-Compatible Respiratory Phantom}{CAD, Instrumentation, Validation}
Developed a programmable motion phantom for validating image-registration algorithms; characterized displacement accuracy to within $\pm$0.6 mm over clinically relevant breathing profiles.

% ---------- Education ----------
\section*{EDUCATION}

\begin{tabularx}{\textwidth}{@{}Xr@{}}
  \textbf{M.S. Biomedical Engineering}, Northeastern University & \textbf{2019} \\
  Concentration: Biomedical Devices and Bioinstrumentation; GPA: 3.86/4.00 & Boston, MA \\
  \textbf{B.S. Biomedical Engineering}, University of Connecticut & \textbf{2017} \\
  Minor: Electrical Engineering; \textit{magna cum laude} & Storrs, CT
\end{tabularx}

% ---------- Skills ----------
\section*{TECHNICAL SKILLS}

\small
\skillrow{Medical-device development}{Design controls, V\&V, requirements traceability, design transfer, CAPA, DHF, usability engineering, ISO 13485, ISO 14971, IEC 60601, FDA 21 CFR 820}
\skillrow{Engineering and analysis}{Python, MATLAB, LabVIEW, Minitab, JMP, statistical analysis, DOE, Gage R\&R, signal processing, test automation}
\skillrow{Design and prototyping}{SolidWorks, Altium Designer, GD\&T, tolerance analysis, PCB bring-up, sensors, microcontrollers, 3D printing, machining}

\end{document}
